% ---
% id: receiver/rem:konig
% kind: rem
% title: "The single-mode normal-ordering identity in the literature"
% summary: >-
%   The single-mode step of the normal-ordered form is the published identity
%   that the normally ordered exponential of z times the number operator is
%   (1 + z) to the number operator, and its λ = 1 case is the perfect on/off
%   detector. Both citations corroborate; neither is used.
% prominence: standard
% status: accepted
% uses: [prop:normord, lem:fock-functor]
% mentions: [thm:click-povm]
% assumes: []
% refs:
%   - "konig2021newton, Eq. (3.23)"
%   - "sperling2012, Eq. (6)"
% provenance: newly-derived
% lean:
% ---

\begin{remark}\label{rem:konig}
Equation \eqref{eq:singlemode} is the single-mode statement
$\normord{e^{z\hat n}}=(1+z)^{\hat n}$ at $z=-\lambda_j$, where
$\hat n=a^\dagger a$. That operator identity is
\cite[Eq.~(3.23)]{konig2021newton}, which reads
$G_f(z;\hat n)=(1+z)^{\hat n}=\normord{e^{z\hat n}}$ and is proved there
without any assumption on the state of the field. Its $\lambda=1$ case,
$\normord{e^{-\hat n}}=\dyad{\vac}$, is the no-click element of a perfect
on/off detector, recorded as \cite[Eq.~(6)]{sperling2012} in the form
$\hat\pi_0=\dyad{0}=\normord{e^{-\hat n}}$, $\hat\pi_1=\normord{\hat 1-e^{-\hat
n}}$. Both are corroboration: \eqref{eq:singlemode} is proved in
Proposition~\ref{prop:normord}, and Theorem~\ref{thm:click-povm} derives the
no-click element from the projector picture of
Lemma~\ref{lem:fock-functor}\,(iv) without using normal ordering at all.
\end{remark}
